% ============================================================
% NLP-P5 · "In RAG We Trust?" — Project working notes
% Course: Natural Language Processing (Prof. A. Ferrara), UniMi
% ============================================================
\documentclass[11pt,a4paper]{report}

% ---------- 1. Fonts and encoding ----------
\usepackage[T1]{fontenc}
\usepackage[utf8]{inputenc}
\usepackage{lmodern}

% ---------- 2. Page layout ----------
\usepackage[margin=2.5cm]{geometry}

% ---------- 3. Mathematics ----------
\usepackage{amsmath, amssymb}

% ---------- 4. Tables ----------
\usepackage{booktabs}
\usepackage{array}
\usepackage{tabularx}
\usepackage{longtable}

% ---------- 5. Lists ----------
\usepackage{enumitem}

% ---------- 6. Colours and boxes ----------
\usepackage{xcolor}
\usepackage[most]{tcolorbox}

\newtcolorbox{definitionbox}[1]{%
  colback=blue!4, colframe=blue!50!black,
  fonttitle=\bfseries, title={Definition: #1}}

\newtcolorbox{examplebox}{%
  colback=gray!6, colframe=gray!60!black,
  fonttitle=\bfseries, title={Example}}

\newtcolorbox{notebox}{%
  colback=yellow!8, colframe=orange!70!black,
  fonttitle=\bfseries, title={Note}}

\newtcolorbox{decisionbox}[1]{%
  colback=green!5, colframe=green!40!black,
  fonttitle=\bfseries, title={Decision #1}}

% ---------- 7. Citations ----------
\usepackage[round,authoryear]{natbib}

% ---------- 8. Hyperlinks (always load last) ----------
\usepackage[colorlinks=true,
            linkcolor=blue!50!black,
            citecolor=green!40!black,
            urlcolor=blue!70!black]{hyperref}

% ---------- 9. Title information ----------
\title{\textbf{In RAG We Trust?}\\[4pt]
       \large Project Working Notes --- NLP Final Project (P5)}
\author{Angela\\ Università degli Studi di Milano}
\date{\today}

% ============================================================
\begin{document}

\pagenumbering{alph}    % title page gets letter numbering (avoids a duplicate "page 1")
\maketitle
\pagenumbering{arabic}  % normal numbering 1, 2, 3 ... from here on
\tableofcontents

\chapter{Project Overview}\label{ch:overview}

\chapter{Literature Audit (A1)}\label{ch:literature}
\section{Scope of the audit}\label{sec:audit-scope}

The project specification for P5 lists three references
\citep{lewis2020rag,zhou2024trustworthiness,singal2024evidence} and states
that they are a starting guideline. We add the two empirical studies closest
to our experiment \citep{jin2024tugofwar,wu2024clasheval}: both measure how
LLMs behave when retrieved evidence is false or conflicting.
Table~\ref{tab:audit} summarises each work.

\begin{notebox}
This area is well studied. Our study does not claim a new finding; its
contribution is a careful, controlled measurement and an honest discussion
of the results.
\end{notebox}

\section{Annotated table}\label{sec:audit-table}

{\small
\renewcommand{\arraystretch}{1.25}
\begin{longtable}{
  >{\raggedright\arraybackslash}p{2.2cm}
  >{\raggedright\arraybackslash}p{3.9cm}
  >{\raggedright\arraybackslash}p{3.9cm}
  >{\raggedright\arraybackslash}p{4.2cm}}
\caption{Annotated literature audit.}\label{tab:audit}\\
\toprule
\textbf{Reference} & \textbf{What they did} & \textbf{What they found} & \textbf{Relation to our study} \\
\midrule
\endfirsthead
\toprule
\textbf{Reference} & \textbf{What they did} & \textbf{What they found} & \textbf{Relation to our study} \\
\midrule
\endhead
\bottomrule
\endfoot

\citet{lewis2020rag}
& Combined a pre-trained seq2seq generator with a dense vector index of
  Wikipedia, searched by a neural retriever; trained the two together.
& State of the art on three open-domain QA tasks; more specific, diverse
  and factual generation than a model without retrieval.
& Origin of RAG. It assumes the retrieved documents are correct; we test
  what happens when they are not. We use the modern form (frozen
  instruction-tuned LLM + separate retriever, no training). \\
\midrule

\citet{zhou2024trustworthiness}
& Survey. Organised RAG trustworthiness into six dimensions: factuality,
  robustness, fairness, transparency, accountability, privacy.
& Most RAG research targets accuracy and efficiency; trustworthiness is
  under-explored. Inappropriate or poorly used retrieval can produce
  undesirable output.
& Gives our vocabulary: the study sits in the \emph{robustness} and
  \emph{factuality} dimensions. \\
\midrule

\citet{singal2024evidence}
& RAG pipeline that retrieves evidence sentences and gives them, with the
  claim, to an LLM for classification; few-shot in-context learning;
  AVeriTeC dataset.
& AVeriTeC score of 0.33, a 22\% absolute improvement over the baseline.
& An evidence-based system whose output depends on the evidence being
  honest. We stress exactly that dependency. \\
\midrule

\citet{jin2024tugofwar}
& Evaluation framework for knowledge conflicts: (i) internal memory vs.\
  external sources; (ii) truthful vs.\ irrelevant vs.\ misleading evidence.
  Proposed a decoding method (CD2).
& Models follow a ``majority rule'' (trust evidence that appears more
  often) and show confirmation bias; stronger models can stick to wrong
  internal memory.
& Closest prior result. Our accuracy-vs-poison-ratio curve partly
  \emph{replicates} the majority-rule finding. We add poison types,
  meta-prompts and hedging behaviour, on open 4-bit models. \\
\midrule

\citet{wu2024clasheval}
& Benchmark of 1{,}200+ questions in six domains; answers in the documents
  perturbed from subtle to blatant errors; six LLMs tested.
& Models adopted incorrect retrieved content, overriding a correct prior,
  over 60\% of the time; the less realistic the error, the less often it
  was adopted.
& Poison plausibility matters, which motivates several poison types. Also a
  warning: blatant poisons would make models look more robust than they
  are. \\

\end{longtable}
}
\section{Positioning of our study}\label{sec:positioning}

We study how open, instruction-tuned language models behave inside a RAG
pipeline when part of the retrieved evidence is deliberately false.
That language models can be misled by false context, and that they tend to
follow the answer supported by most passages, is already established
\citep{jin2024tugofwar,wu2024clasheval}; we do not claim these findings as new.
Our contribution is a controlled measurement in a single, fully reproducible
setting: we vary the share of poisoned passages, the kind of poison, and the
instructions given to the model, and we measure not only whether the final
answer is correct but also whether the model notices the conflict, hedges,
or abstains.
Two limits frame the results from the start. First, we test small open models
in 4-bit precision on a single GPU, so the findings may not transfer to larger
or proprietary models. Second, our poisoned passages are synthetic; real-world
contradictions are often subtler, so our numbers describe a controlled stress
test rather than everyday conditions.
Within these limits, the study asks a practical question: do simple
prompt-level instructions give a measurable protection against poisoned
retrieval, at a cost that a student-level compute budget can afford?


\chapter{Research Questions and Formal Problem 
(A2)}\label{ch:rq}
\section{Definitions}\label{sec:definitions}

This section fixes the vocabulary of the study. Every later research
question, metric and table uses these terms with exactly this meaning.
Concrete values (datasets, number of passages, which models) are chosen in
the experimental design (Chapter~\ref{ch:design}).

\begin{definitionbox}{Question and gold answer}
A \emph{question} $q$ is a factual question taken from a question-answering
dataset. Each question has one \emph{gold answer} $a^{*}$: the answer the
dataset marks as correct.
\end{definitionbox}

\begin{definitionbox}{Passage and context}
A \emph{passage} $p$ is a short piece of text (a few sentences). The
\emph{context} of question $q$ is the ordered list of $k$ passages shown to
the model together with the question:
\[
  C(q) = (p_1, p_2, \dots, p_k).
\]
The order matters: the same passages in a different order form a different
context.
\end{definitionbox}

\begin{definitionbox}{Clean passage}
A passage is \emph{clean} if it is consistent with the gold answer $a^{*}$.
\end{definitionbox}

\begin{definitionbox}{Poisoned passage and target answer}
A passage is \emph{poisoned} if it was written on purpose to support a
specific false answer $\tilde{a} \neq a^{*}$, called the \emph{target
answer}. All poisoned passages for the same question support the
\emph{same} target answer $\tilde{a}$.
\end{definitionbox}

\begin{definitionbox}{Poison ratio}
If $m$ of the $k$ passages in $C(q)$ are poisoned, the \emph{poison ratio}
is
\[
  \rho = \frac{m}{k}, \qquad 0 \le \rho \le 1 .
\]
$\rho = 0$ is a fully clean context; $\rho = 1$ is a fully poisoned one.
\end{definitionbox}

\begin{definitionbox}{Poison type}
The \emph{poison type} $\tau$ is the method used to build the poisoned
passages (for example, swapping the correct entity for another one, or
adding a false claim of authority). $\mathcal{T}$ denotes the set of poison
types studied.
\end{definitionbox}

\begin{definitionbox}{Prompt condition}
The \emph{prompt condition} $\pi$ is the instruction given to the model
together with the question and the context (for example, a plain
instruction to answer, or an instruction to check the passages for
consistency first). $\Pi$ denotes the set of prompt conditions studied.
\end{definitionbox}

\begin{definitionbox}{Model and response}
A \emph{model} $M$ is a language model used without any training. Given a
prompt condition, a question and a context, it produces a free-text
\emph{response}
\[
  r = M\big(\pi, q, C(q)\big).
\]
From $r$ we extract the model's \emph{final answer} $\hat{a}$; if the model
gives no answer, $\hat{a} = \varnothing$.
\end{definitionbox}

\begin{definitionbox}{Outcome}
Each response receives exactly one \emph{outcome} $o \in
\{\textsc{Correct}, \textsc{Poisoned}, \textsc{OtherWrong},
\textsc{Abstain}\}$:
\begin{itemize}[nosep]
  \item \textsc{Correct}: $\hat{a}$ matches the gold answer $a^{*}$;
  \item \textsc{Poisoned}: $\hat{a}$ matches the target answer $\tilde{a}$;
  \item \textsc{OtherWrong}: $\hat{a}$ is given but matches neither;
  \item \textsc{Abstain}: no answer is given ($\hat{a} = \varnothing$).
\end{itemize}
\end{definitionbox}

\begin{definitionbox}{Conflict flag}
Independently of the outcome, a response has a \emph{conflict flag}
$f \in \{0, 1\}$: $f = 1$ if the response explicitly states that the
passages disagree or that some of them may be unreliable, and $f = 0$
otherwise.
\end{definitionbox}

\begin{notebox}
The outcome and the conflict flag are kept separate on purpose: a model can
notice the conflict \emph{and still} answer correctly, wrongly, or not at
all. The exact rules for deciding when two answers ``match'' are fixed in
Chapter~\ref{ch:design}.
\end{notebox}
\section{Formal problem definition}\label{sec:formal}

\paragraph{Data.}
Let $\mathcal{Q} = \{(q_i, a^{*}_i)\}_{i=1}^{N}$ be a set of $N$ questions
with their gold answers. For each question $q_i$ we have:
\begin{itemize}[nosep]
  \item a pool of clean passages $P^{+}(q_i)$;
  \item one target answer $\tilde{a}_i \neq a^{*}_i$;
  \item for each poison type $\tau \in \mathcal{T}$, a pool of poisoned
        passages $P^{-}_{\tau}(q_i)$, all supporting $\tilde{a}_i$.
\end{itemize}

\paragraph{Context construction.}
Given a poison ratio $\rho$ and a poison type $\tau$, the context of $q_i$
is built by taking $m = \rho k$ passages from $P^{-}_{\tau}(q_i)$ and
$k - m$ passages from $P^{+}(q_i)$, and placing them in a random order
$\sigma_i$:
\[
  C_{\rho,\tau}(q_i) = \sigma_i\big(p^{-}_{1}, \dots, p^{-}_{m},\;
                                   p^{+}_{1}, \dots, p^{+}_{k-m}\big).
\]

\paragraph{Experimental condition.}
An \emph{experimental condition} is a tuple
\[
  \theta = (M, \pi, \tau, \rho) \in
  \mathcal{M} \times \Pi \times \mathcal{T} \times \mathcal{R},
\]
where $\mathcal{M}$ is the set of models and $\mathcal{R}$ the set of
poison ratios studied. Under condition $\theta$, question $q_i$ produces
\[
  r_i(\theta) = M\big(\pi, q_i, C_{\rho,\tau}(q_i)\big),
\]
with outcome $o_i(\theta)$ and conflict flag $f_i(\theta)$ as defined in
Section~\ref{sec:definitions}.

\paragraph{Quantities of interest.}
For each condition $\theta$ we estimate four rates over the $N$ questions,
where $\mathbb{1}[\cdot]$ is $1$ if the statement inside is true and $0$
otherwise:
\begin{align*}
  \mathrm{Acc}(\theta)     &= \tfrac{1}{N}\textstyle\sum_{i=1}^{N}
                              \mathbb{1}\big[o_i(\theta)=\textsc{Correct}\big],\\
  \mathrm{ASR}(\theta)     &= \tfrac{1}{N}\textstyle\sum_{i=1}^{N}
                              \mathbb{1}\big[o_i(\theta)=\textsc{Poisoned}\big],\\
  \mathrm{Abst}(\theta)    &= \tfrac{1}{N}\textstyle\sum_{i=1}^{N}
                              \mathbb{1}\big[o_i(\theta)=\textsc{Abstain}\big],\\
  \mathrm{Flag}(\theta)    &= \tfrac{1}{N}\textstyle\sum_{i=1}^{N} f_i(\theta).
\end{align*}
$\mathrm{Acc}$ is the accuracy; $\mathrm{ASR}$ is the \emph{attack success
rate}, the share of answers in which the poison succeeded; $\mathrm{Abst}$
is the abstention rate; $\mathrm{Flag}$ is the conflict-flag rate.

\paragraph{Problem.}
The study estimates how these rates change as one factor of $\theta$ varies
while the others are held fixed: as a function of $\rho$ (robustness
curve), between prompt conditions $\pi$ (effect of meta-prompts), between
poison types $\tau$, and between models $M$. Because the same $N$ questions
are used in every condition, all comparisons are \emph{paired}: each
question serves as its own control.

\begin{notebox}
Statistical estimators, confidence intervals, self-consistency and the
remaining course metrics are defined in Chapter~\ref{ch:design}. The value
of $k$ and the sets $\mathcal{M}$, $\Pi$, $\mathcal{T}$, $\mathcal{R}$ are
fixed there; they must make $m = \rho k$ a whole number.
\end{notebox}

\section{Hypotheses}\label{sec:hypotheses}

The hypotheses below were written \emph{before} any experiment was run.
Where a hypothesis follows from prior work, the source is cited; where it
is our own expectation, this is stated. A result that contradicts a
hypothesis is reported and discussed, not hidden.

\begin{description}[leftmargin=0pt]
  \item[H1 (for RQ1).] As the poison ratio $\rho$ grows, $\mathrm{Acc}$
    decreases and $\mathrm{ASR}$ increases. The change is steepest around
    $\rho = 0.5$, where the poisoned passages become the majority,
    following the ``majority rule'' reported by \citet{jin2024tugofwar}.
    At $\rho = 1$, $\mathrm{ASR}$ is high even for questions the model
    could answer correctly without context, as models often override a
    correct prior with false context \citep{wu2024clasheval}.

  \item[H2 (for RQ2).] \emph{Our expectation.} Meta-prompts lower
    $\mathrm{ASR}$ mainly at intermediate ratios ($0 < \rho < 1$), where
    clean passages are present to compare against; at $\rho = 1$ they
    help little, because no clean passage contradicts the poison. On clean
    contexts ($\rho = 0$) they cause a small loss of accuracy and a higher
    abstention rate.

  \item[H3 (for RQ3).] \emph{Our expectation.} $\mathrm{Flag}$ is low at
    $\rho = 0$, highest at intermediate ratios, and low again at
    $\rho = 1$, because only mixed contexts contain passages that disagree
    with each other. Responses with $f = 1$ show a lower $\mathrm{ASR}$
    than responses with $f = 0$ under the same condition.

  \item[H4 (for RQ4).] Poison types that look more plausible or more
    authoritative reach a higher $\mathrm{ASR}$, consistent with the
    finding that less realistic false content is adopted less often
    \citep{wu2024clasheval}. Models differ in robustness, but we make no
    directional prediction about which model is more robust: prior work
    reports that stronger models can also cling to wrong internal
    knowledge \citep{jin2024tugofwar}.
\end{description}

\begin{notebox}
H3 reveals a property of the design: at $\rho = 1$ there is no conflict
\emph{between passages}, only between the passages and what the model may
already know. A low $\mathrm{Flag}$ at $\rho = 1$ is therefore expected and
does not by itself mean the model is careless.
\end{notebox}


\section{Measurable objectives}\label{sec:objectives}

Each research question is answered by a specific set of numbers. Every
estimate is reported with a 95\% confidence interval; the estimation and
testing procedures are fixed in Chapter~\ref{ch:design}.

\begin{description}[leftmargin=0pt]
  \item[O1 (RQ1).] For every model, prompt condition and poison type,
    estimate $\mathrm{Acc}$ and $\mathrm{ASR}$ at each poison ratio
    $\rho \in \mathcal{R}$, giving one robustness curve per setting.
    Report the smallest ratio at which $\mathrm{ASR}$ exceeds
    $\mathrm{Acc}$ (the \emph{crossover point}).

  \item[O2 (RQ2).] For each meta-prompt $\pi$ and each $\rho$, estimate
    the paired difference from the plain prompt $\pi_0$:
    \[
      \Delta\mathrm{ASR}_{\pi}(\rho) =
      \mathrm{ASR}(\pi, \rho) - \mathrm{ASR}(\pi_0, \rho),
    \]
    with model and poison type held fixed, and test whether it differs
    from zero. On clean contexts ($\rho = 0$), estimate the corresponding
    differences in $\mathrm{Acc}$ and $\mathrm{Abst}$ (the cost of the
    meta-prompt).

  \item[O3 (RQ3).] Estimate $\mathrm{Flag}$ at each $\rho$. Within each
    condition, compare the attack success rate of flagged and unflagged
    responses,
    \[
      \mathrm{ASR}(\theta \mid f = 1) \quad\text{vs.}\quad
      \mathrm{ASR}(\theta \mid f = 0),
    \]
    and test whether the difference is non-zero.

  \item[O4 (RQ4).] Compare the quantities of O1--O3 across poison types
    $\tau$ and models $M$, and test whether the effects of $\rho$ and of
    the prompt condition depend on $\tau$ or $M$.

  \item[O5 (course metrics).] Report the hallucination rate and the
    self-consistency of answers, as required by the project
    specification; both are defined in Chapter~\ref{ch:design}.
\end{description}







\chapter{Experimental Design (A3)}\label{ch:design}

\chapter{Environment and Repository (A4)}\label{ch:repo}

\chapter{Decision Log}\label{ch:decisions}
Each decision records what was chosen, why, which alternatives were
rejected, and a one-sentence justification for the oral discussion.

\section{Block A}\label{sec:decisions-a}

\begin{decisionbox}{D1 --- Repository created at the end of A1}
\begin{description}[nosep, leftmargin=0pt]
  \item[Decision.] Create the GitHub repository at the end of A1 with only a
    README and a \texttt{docs/} folder; A4 builds the full package layout
    inside the same repository.
  \item[Why.] Every phase must end with a push, and the commit history then
    records how the methodology was designed.
  \item[Alternatives.] Push everything only in A4 (breaks the push-per-phase
    rule, risks losing work); build the full layout now (does A4's work out
    of order).
  \item[Interview line.] ``I versioned the research design from day one, so
    the repository shows how the methodology was decided.''
\end{description}
\end{decisionbox}

\begin{decisionbox}{D2 --- Cite version 1 of the Zhou et al.\ survey}
\begin{description}[nosep, leftmargin=0pt]
  \item[Decision.] Cite arXiv:2409.10102 version 1 (September 2024).
  \item[Why.] It is the version listed in the course specification; a 2026
    revision changed the author list and added a benchmark.
  \item[Alternatives.] Cite only the latest version (author list would not
    match the course reference).
  \item[Interview line.] ``I cite the version assigned in the course and I am
    aware of the later revision.''
\end{description}
\end{decisionbox}

\begin{decisionbox}{D3 --- Spelling ``Singal''}
\begin{description}[nosep, leftmargin=0pt]
  \item[Decision.] Use the spelling ``Singal'' from the ACL Anthology record.
  \item[Why.] It is the official published record and matches the course
    specification; the authors' repository spells it ``Singhal''.
  \item[Alternatives.] Use the repository spelling (not the official record).
  \item[Interview line.] ``I follow the official proceedings record.''
\end{description}
\end{decisionbox}

\begin{decisionbox}{D4 --- Search for related work beyond the course list}
\begin{description}[nosep, leftmargin=0pt]
  \item[Decision.] Search for closely related work missing from the plan;
    five verified papers were identified. \emph{Partly superseded by D7:}
    only two are kept.
  \item[Why.] The papers closest to our experiment protect the paper from
    presenting known results as new.
  \item[Alternatives.] Rely only on the course list (leaves the closest
    empirical work uncited).
  \item[Interview line.] ``I checked for the closest prior studies before
    defining my research questions.''
\end{description}
\end{decisionbox}

\begin{decisionbox}{D5 --- Three-level reading plan}
\begin{description}[nosep, leftmargin=0pt]
  \item[Decision.] Read the three course references; know Jin et al.\ and
    ClashEval from verified summaries; read further papers only if time
    allows.
  \item[Why.] The project must be completed in 16 days; the course grades the
    study and its critical discussion, not the number of papers read.
  \item[Alternatives.] Read every candidate paper (not feasible in the time
    available).
  \item[Interview line.] ``I studied the course references in depth and
    positioned my work against the closest empirical studies.''
\end{description}
\end{decisionbox}

\begin{decisionbox}{D6 --- LaTeX working notes, separate from the paper}
\begin{description}[nosep, leftmargin=0pt]
  \item[Decision.] Keep these working notes in the \texttt{report} class; the
    final 4--8 page paper is written in the Springer template in Block G.
  \item[Why.] The paper template and page limit are fixed; the notes collect
    the full design and are condensed into the paper later.
  \item[Alternatives.] Start the Springer paper now (wrong place for notes,
    belongs to Block G); Markdown notes (LaTeX was requested).
  \item[Interview line.] ``I kept a versioned research notebook during design
    and condensed it into the paper.''
\end{description}
\end{decisionbox}

\begin{decisionbox}{D7 --- Five references (Option A)}
\begin{description}[nosep, leftmargin=0pt]
  \item[Decision.] Cite the three course references plus Jin et al.\ (2024)
    and ClashEval; the eight candidates from the plan are not cited.
  \item[Why.] The course PDF lists exactly three references as a starting
    guideline and requires only a short literature discussion; the two added
    studies are the closest to our experiment. Every entry was checked
    against an official source.
  \item[Alternatives.] Also cite the eight candidates (richer, but costly in
    time and space in a 4--8 page paper).
  \item[Interview line.] ``I built on the course references and positioned my
    results against the two closest empirical studies.''
\end{description}
\end{decisionbox}

\begin{decisionbox}{D8 --- ``Relation to our study'' instead of ``Gap''}
\begin{description}[nosep, leftmargin=0pt]
  \item[Decision.] The last column of Table~\ref{tab:audit} describes how
    each work relates to ours, not what it ``missed''.
  \item[Why.] Claiming a gap requires full knowledge of each paper's setup;
    an unverified gap claim could be wrong.
  \item[Alternatives.] A ``Gap'' column (pushes toward inflated claims).
  \item[Interview line.] ``I described the relation to each work without
    claiming gaps I could not verify.''
\end{description}
\end{decisionbox}

\begin{decisionbox}{D9 --- No novelty claim; limits stated up front}
\begin{description}[nosep, leftmargin=0pt]
  \item[Decision.] The positioning (Section~\ref{sec:positioning}) contains no
    novelty claim and states the two main limits: small 4-bit models and
    synthetic poisons.
  \item[Why.] The area is well studied; naming weaknesses before a reviewer
    does is part of an honest discussion.
  \item[Alternatives.] Frame the study as novel (easy to attack and untrue).
  \item[Interview line.] ``My contribution is rigour and honest discussion in
    a controlled, reproducible setting.''
\end{description}
\end{decisionbox}



\begin{decisionbox}{D10 --- Develop on GitLab, deliver on GitHub}
\begin{description}[nosep, leftmargin=0pt]
  \item[Decision.] Develop in a private GitLab repository; upload the full
    history to a public GitHub repository for submission.
  \item[Why.] The course requires a GitHub link; a regional restriction
    blocks private GitHub repositories.
  \item[Alternatives.] Public GitHub from the start (unfinished work
    public); GitLab only (does not meet the written requirement).
  \item[Interview line.] ``I developed privately and published the complete
    history on GitHub for submission.''
\end{description}
\end{decisionbox}

\begin{decisionbox}{D11 --- VS Code without AI code completion}
\begin{description}[nosep, leftmargin=0pt]
  \item[Decision.] Use Visual Studio Code with no built-in AI completion.
  \item[Why.] Free, native on Apple Silicon, handles Python, notebooks,
    terminal and git; keeps every AI contribution traceable in the log.
  \item[Alternatives.] PyCharm (heavier, some features paid); AI-first
    editors (untraceable AI contributions).
  \item[Interview line.] ``Every AI contribution to my code is in my usage
    log, because my editor adds none silently.''
\end{description}
\end{decisionbox}

\begin{decisionbox}{D12 --- Git identity}
\begin{description}[nosep, leftmargin=0pt]
  \item[Decision.] Commits are signed as Angela Aldaoud with a personal
    email address, chosen by the student.
  \item[Why.] One identity for all commits; the address is linked to the
    hosting accounts.
  \item[Alternatives.] University address (initially proposed).
  \item[Interview line.] ``All commits carry one consistent author
    identity.''
\end{description}
\end{decisionbox}

\begin{decisionbox}{D13 --- One target answer per question}
\begin{description}[nosep, leftmargin=0pt]
  \item[Decision.] All poisoned passages for a question support the same
    false answer $\tilde{a}$.
  \item[Why.] Makes poison adoption measurable and separable from ordinary
    errors; $\rho$ then measures the strength of one coherent attack.
  \item[Alternatives.] A different false answer per passage (models noise,
    but blurs the measurement).
  \item[Interview line.] ``One target answer per question lets me tell
    poison adoption apart from other mistakes.''
\end{description}
\end{decisionbox}

\begin{decisionbox}{D14 --- Four outcomes}
\begin{description}[nosep, leftmargin=0pt]
  \item[Decision.] Correct, Poisoned, OtherWrong, Abstain.
  \item[Why.] ``Wrong'' hides two different failures: adopting the poison
    and failing for other reasons.
  \item[Alternatives.] Right/wrong only (cannot show whether the poison
    worked).
  \item[Interview line.] ``I separate being fooled by the poison from
    other errors.''
\end{description}
\end{decisionbox}

\begin{decisionbox}{D15 --- Conflict flag separate from the outcome}
\begin{description}[nosep, leftmargin=0pt]
  \item[Decision.] The conflict flag $f$ is recorded independently of the
    outcome.
  \item[Why.] A model can notice a conflict and still answer correctly,
    wrongly or not at all; merging the two loses that information.
  \item[Alternatives.] A combined ``abstained/flagged'' category (from the
    original plan).
  \item[Interview line.] ``I measure whether the model is right and whether
    it notices the conflict, as two separate things.''
\end{description}
\end{decisionbox}

\begin{decisionbox}{D16 --- Random passage order}
\begin{description}[nosep, leftmargin=0pt]
  \item[Decision.] Passages in each context are placed in a random order.
  \item[Why.] Models use passages differently depending on position; a
    fixed position would confound position with poison.
  \item[Alternatives.] Fixed position (biased); position as an extra factor
    (a separate question).
  \item[Interview line.] ``I randomised passage order so position effects
    cannot pass as poison effects.''
\end{description}
\end{decisionbox}

\begin{decisionbox}{D17 --- Paired design}
\begin{description}[nosep, leftmargin=0pt]
  \item[Decision.] The same questions are used in every condition.
  \item[Why.] Each question is its own control: comparisons are more
    precise and paired tests become possible.
  \item[Alternatives.] Different questions per condition (needs more data,
    weaker comparisons).
  \item[Interview line.] ``Each question serves as its own control.''
\end{description}
\end{decisionbox}

\begin{decisionbox}{D18 --- Attack success rate as a main metric}
\begin{description}[nosep, leftmargin=0pt]
  \item[Decision.] Report ASR next to accuracy.
  \item[Why.] Accuracy falls for two reasons; ASR isolates the one P5 is
    about. It is the standard name in the poisoning literature.
  \item[Alternatives.] Accuracy only.
  \item[Interview line.] ``ASR tells me how often the poison actually
    succeeded.''
\end{description}
\end{decisionbox}

\begin{decisionbox}{D19 (revised) --- All four research questions run in full}
\begin{description}[nosep, leftmargin=0pt]
  \item[Decision.] RQ1--RQ4 are all answered on the full experimental grid;
    RQ4 is secondary only in the space it receives in the discussion.
  \item[Why.] Student's instruction: the grid is sized to the statistical
    needs of the questions, and compute limits are solved with resources,
    not by cutting the design. \emph{Replaces the earlier version, in which
    RQ4 was cuttable.}
  \item[Alternatives.] Make RQ4 cuttable (rejected by the student).
  \item[Interview line.] ``I sized the experiment to the questions, then
    found the compute to run it.''
\end{description}
\end{decisionbox}

\begin{decisionbox}{D20 --- RQ2 measures the cost of meta-prompts}
\begin{description}[nosep, leftmargin=0pt]
  \item[Decision.] RQ2 includes accuracy on clean contexts and extra
    abstentions.
  \item[Why.] A prompt that lowers ASR by refusing to answer is not an
    improvement.
  \item[Alternatives.] Measure only the reduction in ASR.
  \item[Interview line.] ``I measured what a protective prompt costs, not
    only what it gains.''
\end{description}
\end{decisionbox}

\begin{decisionbox}{D21 --- RQ3 is an association}
\begin{description}[nosep, leftmargin=0pt]
  \item[Decision.] RQ3 is stated and analysed as an association, not a
    causal effect.
  \item[Why.] The conflict flag is an output of the model, not a variable
    the experimenter controls.
  \item[Alternatives.] A causal claim (not supported by the design).
  \item[Interview line.] ``RQ3 is correlational, because I do not
    manipulate the flag.''
\end{description}
\end{decisionbox}

\begin{decisionbox}{D22 --- Colab Pro for one month}
\begin{description}[nosep, leftmargin=0pt]
  \item[Decision.] Subscribe to Colab Pro for one month to run the full grid.
  \item[Why.] The grid is not reduced; the free T4 alone may not fit it
    before 15 October. Colab Pro uses compute units, so the A3 budget must
    check that the grid fits.
  \item[Alternatives.] University GPU group (granted for the thesis; not
    used for this project); free T4 only (may not fit).
  \item[Interview line.] ``I paid for the compute the design needed rather
    than shrinking the design.''
\end{description}
\end{decisionbox}

\begin{decisionbox}{D23 (revised) --- GitHub upload at submission}
\begin{description}[nosep, leftmargin=0pt]
  \item[Decision.] Upload to GitHub on submission day; the repository link
    is already known and can be given in the booking form.
  \item[Why.] Configuring GitHub access now costs time with no benefit
    before submission. \emph{Replaces the earlier plan of publishing on
    GitHub immediately.}
  \item[Alternatives.] Publish on GitHub now.
  \item[Interview line.] ``The public repository contains the full
    development history.''
\end{description}
\end{decisionbox}

\begin{decisionbox}{D24 --- Hypotheses declare their source}
\begin{description}[nosep, leftmargin=0pt]
  \item[Decision.] Each hypothesis is marked as cited or as our own
    expectation, and written before any experiment.
  \item[Why.] Results can then be read honestly against prior work and
    against our own intuition.
  \item[Alternatives.] Unsourced hypotheses, or hypotheses written after
    the results.
  \item[Interview line.] ``I wrote my hypotheses before running anything
    and marked which came from prior work.''
\end{description}
\end{decisionbox}

\begin{decisionbox}{D25 --- No directional prediction between models}
\begin{description}[nosep, leftmargin=0pt]
  \item[Decision.] H4 predicts that models differ, not which one is more
    robust.
  \item[Why.] Prior work gives no clear direction; a guessed direction
    would not be honest. Leads to two-sided tests.
  \item[Alternatives.] Predict that the larger model is more robust.
  \item[Interview line.] ``I had no evidence for a direction, so I tested
    two-sided.''
\end{description}
\end{decisionbox}

\begin{decisionbox}{D26 --- 95\% confidence intervals everywhere}
\begin{description}[nosep, leftmargin=0pt]
  \item[Decision.] Every estimate is reported with a 95\% confidence
    interval.
  \item[Why.] A rate without its uncertainty cannot support a claim; the
    sample size is chosen so that the intervals are tight.
  \item[Alternatives.] Point estimates only.
  \item[Interview line.] ``Every number in my paper has a confidence
    interval.''
\end{description}
\end{decisionbox}

\begin{decisionbox}{D27 --- Crossover point}
\begin{description}[nosep, leftmargin=0pt]
  \item[Decision.] Summarise each robustness curve by the smallest $\rho$
    at which ASR exceeds accuracy; ``not reached'' if it never does.
  \item[Why.] One comparable number per setting for tables and slides.
  \item[Alternatives.] Compare full curves only.
  \item[Interview line.] ``The crossover point says when the model starts
    believing the lie more than the truth.''
\end{description}
\end{decisionbox}

\begin{decisionbox}{D28 --- Course metrics kept explicitly}
\begin{description}[nosep, leftmargin=0pt]
  \item[Decision.] Hallucination rate and self-consistency are reported
    (objective O5), defined in the experimental design.
  \item[Why.] The project specification lists them as evaluation metrics.
  \item[Alternatives.] Leave them out silently.
  \item[Interview line.] ``I report every metric the specification asks
    for.''
\end{description}
\end{decisionbox}



\chapter{AI-Usage Log}\label{ch:ailog}
The course requires a disclaimer stating which AI models were used, for what
purposes, and to what extent their outputs were modified or verified. This
chapter is the running record from which that disclaimer will be written.

\begin{notebox}
Model used in Block A: Claude Opus 5.5 (Anthropic), in a chat interface with
web search. All research decisions were taken by the student; every
bibliographic fact was checked against an official source (NeurIPS, AAAI,
ACL Anthology, arXiv) or the course PDF.
\end{notebox}

\section{Block A}\label{sec:ailog-a}

{\small
\renewcommand{\arraystretch}{1.25}
\begin{longtable}{
  >{\raggedright\arraybackslash}p{0.6cm}
  >{\raggedright\arraybackslash}p{4.4cm}
  >{\raggedright\arraybackslash}p{4.4cm}
  >{\raggedright\arraybackslash}p{4.8cm}}
\caption{AI usage during Block A.}\label{tab:ailog-a}\\
\toprule
\textbf{\#} & \textbf{Purpose} & \textbf{AI output} & \textbf{Verification / modification} \\
\midrule
\endfirsthead
\toprule
\textbf{\#} & \textbf{Purpose} & \textbf{AI output} & \textbf{Verification / modification} \\
\midrule
\endhead
\bottomrule
\endfoot
1 & Recover the project plan from the organizer chat
  & Summary of the plan and of Block A
  & Confirmed by the student against her own chat \\
2 & Plan phase A1 and the working rules
  & Step list; rules for explaining details and justifying decisions
  & Reviewed and accepted by the student; rules added at her request \\
3 & Verify the three course references
  & Citations and plain-language summaries
  & Checked against NeurIPS, arXiv and ACL Anthology records \\
4 & Check the literature requirement
  & Answer based on the requirements extracted earlier
  & Later confirmed directly in the course PDF (entry 8) \\
5 & Search for closely related work
  & Five candidate papers with summaries
  & Each checked against its official record; scope reduced by the student (D7) \\
6 & Plan the reading load for 16 days
  & Three-level reading plan
  & Priority set by the student \\
7 & Set up the LaTeX working notes
  & Preamble and chapter skeleton
  & Test-compiled; one warning found and fixed \\
8 & Confirm the P5 references
  & Extract of the P5 reference list
  & Read directly from the course PDF \\
9 & Build the bibliography
  & \texttt{refs.bib} with five entries
  & Fields copied from official records; test-compiled \\
10 & Annotated literature table
  & Chapter~\ref{ch:literature}, Table~\ref{tab:audit}
  & Content from verified sources only; test-compiled \\
11 & Positioning of the study
  & Section~\ref{sec:positioning}
  & No novelty claim; test-compiled \\
12 & Decision log
  & Chapter~\ref{ch:decisions} (D1--D9)
  & Matches the decisions taken in the chat; test-compiled \\
13 & AI-usage log
  & This chapter
  & Test-compiled \\

14 & Set up git, SSH keys and the GitLab repository on the student's Mac
  & Terminal commands and troubleshooting steps
  & Each command run by the student; every output checked before the next step \\
15 & Choose code hosting and compute
  & Options with trade-offs (GitLab/GitHub, Colab Pro)
  & Final choices made by the student \\
16 & A2 definitions
  & Section~\ref{sec:definitions}
  & Test-compiled \\
17 & Formal problem definition
  & Section~\ref{sec:formal}
  & Test-compiled \\
18 & Research questions
  & Section~\ref{sec:rqs}
  & Revised after the student's correction (no reduction of the grid); test-compiled \\
19 & Hypotheses
  & Section~\ref{sec:hypotheses}
  & Cited findings taken only from verified sources; test-compiled \\
20 & Measurable objectives
  & Section~\ref{sec:objectives}
  & Test-compiled \\
21 & Decision and AI-usage logs for A2
  & D10--D28 and entries 14--21
  & Matches the decisions taken in the chat; test-compiled \\
\end{longtable}
}

% ---------- Bibliography (activated once refs.bib exists) ----------
 \bibliographystyle{plainnat}
 \bibliography{refs}

\end{document}